\documentclass[11pt]{article}
\usepackage[margin=1.15in]{geometry}
\usepackage{amsmath,amssymb,amsthm,mathtools}
\usepackage{xcolor}
\usepackage[colorlinks=true,linkcolor=blue!60!black,citecolor=blue!60!black,urlcolor=blue!60!black]{hyperref}

\newtheorem{theorem}{Theorem}[section]
\newtheorem{proposition}[theorem]{Proposition}
\newtheorem{lemma}[theorem]{Lemma}
\newtheorem{corollary}[theorem]{Corollary}
\theoremstyle{definition}
\newtheorem{definition}[theorem]{Definition}
\theoremstyle{remark}
\newtheorem{remark}[theorem]{Remark}

\newcommand{\C}{\mathbb{C}}
\newcommand{\R}{\mathbb{R}}
\newcommand{\Z}{\mathbb{Z}}
\newcommand{\w}{\omega}
\newcommand{\Sp}{\mathrm{Sp}}
\newcommand{\Stab}{\mathrm{Stab}}
\newcommand{\tr}{\operatorname{tr}}
\newcommand{\diag}{\operatorname{diag}}
\newcommand{\dg}{\dagger}
\newcommand{\arccosh}{\operatorname{arccosh}}

\title{The Pais--Uhlenbeck Metric as a Minimum-Distortion Principle,\\
and the Representation Problem for the Fourth-Order Field}
\author{GPT-5.6.sol\thanks{OpenAI model and principal programme author.}
\and Asger Alstrup Palm\thanks{Programme orchestrator and corresponding human contact; Honorary Professor, DTU Compute, Technical University of Denmark. \texttt{asger@area9.dk}.}}
\date{Open-repository preprint, July 2026}

\begin{document}
\maketitle

\begin{abstract}
For the unequal-frequency PT-symmetric Pais--Uhlenbeck oscillator, the
admissible complex-symplectic diagonalizers form a coset
$S_+\,SO(2,\C)^2$ of the unique Hermitian positive-definite diagonalizer
$S_+$.  We prove that $S_+$ minimizes the affine-invariant metric distortion:
for every admissible $S$,
$\|\log(S^{\dg}S)\|_F^2\ge 4r^2$ with
$r=\log\frac{\w_1+\w_2}{\w_1-\w_2}$, with equality exactly on the unitary
stabilizer coset, on which the induced Hilbert-space metric is constantly the
Bender--Mannheim $\eta=e^{-Q}$.  The proof is elementary and exact: at
worst-case stabilizer alignment the log-eigenvalue pair of $S^{\dg}S$ is
$\{\arccosh(\cosh r\cosh b)+a,\ |\arccosh(\cosh r\cosh b)-a|\}$ in the
Gram rapidities $(a,b)$ of the
stabilizer factor.  The general result behind it is an orthogonal Cartan-norm projection
principle, proved in if-and-only-if form, whose principal case is the
$\Sp(2n,\C)$ inter-mode theorem: relative to any symplectic mode splitting,
minimality along the mode-preserving coset holds exactly for inter-mode
diagonalizer directions.  The Pais--Uhlenbeck oscillator is the fully
integrable $n=2$ model: the physical stabilizer, itself not
$\theta$-stable, has canonical compatible hull the mode-preserving
$SL(2,\C)^2$ (totally geodesic $\mathbb H^3\times\mathbb H^3$;
generically, with the alignment stratification made explicit), the
Bender--Mannheim direction is inter-mode, and the closed form above is the
resulting mixed hyperbolic/flat Pythagorean law.  We then quantize the
free fourth-order field modewise and compute the Bender--Mannheim ground
state exactly: it is a normalizable real Gaussian for every mode ---
indeed exactly the Dyson pullback $\rho^{-1}\phi_0$ of the normal-form
vacuum, so the modewise Dyson maps are \emph{pointed} $\eta$-unitaries
and the physical PT completion is globally unitarily equivalent to the
positive normal-form Fock completion.  If instead the two Gaussian vacua
are compared through the identity embedding of the auxiliary canonical
variables and the standard involution, their identity-embedding Gaussian
fidelity tends to the universal constant $\sqrt3/2$ at large momentum,
their auxiliary standard-Fock occupation to $1/3$ per mode pair, and the
two auxiliary standard-CCR product states are disjoint in every spatial
dimension.  The obstruction is thus not physical Hilbert-space
inequivalence but the incompatibility between the identity embedding of
canonical variables and the Dyson-transported physical involution.  All
identities are machine-verified.
\end{abstract}

\section*{AI authorship and computational accountability}
GPT-5.6.sol developed the derivations, computational checks, claim audit, and
manuscript.  Asger Alstrup Palm commissioned and orchestrated the programme
and is the corresponding human contact.  The project is an experiment in
whether a non-specialist human, using AI systems as research collaborators,
can organize falsifiable computer-assisted mathematical physics.  The exact
status of each repository check is stated where it is used, and no check is
promoted beyond its documented scope.  The manuscript remains subject to
expert review.

\section{Introduction}\label{sec:intro}

This paper continues the symplectic analysis of the PT-symmetric
Pais--Uhlenbeck (PU) oscillator begun in \cite{companion}, whose results and
conventions we use freely: $V=(x,y,p,q)^{\top}$, $[V_a,V_b]=iJ_{ab}$,
$H_{\mathrm{PT}}=\tfrac12V^{\top}GV$ with the PU matrix $G$, positive normal
form $G_0$, rapidity
\begin{equation}
r=\log\frac{\w_1+\w_2}{\w_1-\w_2},\qquad \w_1>\w_2>0,
\end{equation}
and, after the canonical normalization $d_xd_y=\gamma\w_1\w_2$, the unique
Hermitian positive-definite diagonalizer
\begin{equation}
S_+=\cosh\tfrac r2\,I+\sinh\tfrac r2\,B\in\Sp(4,\C),
\qquad
S_+^{\top}G'S_+=G_0',\qquad S_+^{\top}JS_+=J,
\end{equation}
with $B$ a fixed Hermitian involution.  The full solution family of the
diagonalization problem is the coset $S_+\,\Stab(J,G_0')$ with
$\Stab(J,G_0')\cong SO(2,\C)\times SO(2,\C)$, and the metaplectic logarithm
of $S_+$ is the Bender--Mannheim generator $Q=\alpha pq+\beta xy$
\cite{BM2008PRL,BM2008PRD,companion}.

Two questions are addressed here.

\emph{First}, uniqueness of the positive diagonalizer does not by itself
distinguish $S_+$ variationally: one may ask whether $S_+$ \emph{minimizes} a
natural functional over all admissible diagonalizers.  We prove that it
minimizes the affine-invariant distortion
$\mathcal F(S)=d(I,S^{\dg}S)^2=\|\log(S^{\dg}S)\|_F^2$
(Theorem~\ref{thm:mindist}), so that the Bender--Mannheim metric is selected
by a minimum-distortion principle: it is the least-distorting complex
canonical transformation converting the PU Hamiltonian to its positive
normal form.  The results of Section~\ref{sec:mindist} form a hierarchy:
\begin{equation*}
\text{normal-coset projection principle (Thm.~\ref{thm:projection})}
\;\Longrightarrow\;
\Sp(2n,\C)\ \text{inter-mode theorem (Thm.~\ref{thm:nmode})}
\;\Longrightarrow\;
\text{exact }\mathbb H^2\times\R\ \text{formula \eqref{eq:xformula}} .
\end{equation*}
The first is classical symmetric-space geometry, stated here in
if-and-only-if form; the second is the principal general result, making the
selection principle a theorem for arbitrary quadratic PT mode systems; the
third is what only the integrable $n=2$ case supplies --- the exact
minimizing value, obtained by the elementary two-invariant derivation that
opens the section.  The physical stabilizer itself fails
$\theta$-stability --- which is why the classical theory does not apply to
it directly --- but its canonical compatible hull is exactly the
mode-preserving $SL(2,\C)^2$ (Proposition~\ref{prop:hull}, stratified over
the alignment locus).

\emph{Second}, for the free fourth-order field
$(\Box+m_1^2)(\Box+m_2^2)\phi=0$ each spatial momentum $k$ contributes a PU
pair with $\w_i(k)=\sqrt{k^2+m_i^2}$, and the modewise construction
applies with rapidity
\begin{equation}\label{eq:rk}
r(k)=\log\frac{\big(\w_1(k)+\w_2(k)\big)^2}{m_1^2-m_2^2}
=2\log|k|+O(1)\qquad(m_1>m_2).
\end{equation}
One might expect from \eqref{eq:rk} that the modewise transformation carries
unboundedly growing squeezing and fails Fock implementability
``dramatically''.  The exact computation shows something sharper and more
structured (Theorems~\ref{thm:groundstate}--\ref{thm:disjoint} and
Proposition~\ref{prop:pointed}): the PT
ground state exists as a normalizable Gaussian for \emph{every} mode, and
it is exactly the Dyson pullback $\psi_0=\rho^{-1}\phi_0$ of the
normal-form vacuum.  Two comparisons must then be sharply distinguished.
\emph{Physical completion}: the modewise Dyson maps are pointed
$\eta$-unitaries, their pointed infinite tensor product exists, and the
physical PT completion is globally unitarily equivalent to the positive
normal-form Fock completion (Proposition~\ref{prop:pointed}).
\emph{Auxiliary identity embedding}: if instead both Gaussian vacua are
regarded as vectors of one standard $L^2$ space, with the auxiliary
canonical variables identified by the identity map and the standard
involution retained, their mismatch saturates at \emph{universal
constants} --- identity-embedding fidelity $\sqrt3/2$ and auxiliary
standard-Fock occupation $1/3$ per mode pair at large $k$ --- and the two
auxiliary standard-CCR product states are disjoint in every spatial
dimension $d\ge1$, with overlap
$|\langle\Omega,\Psi\rangle|^2
=\exp[-(\log\tfrac{2}{\sqrt3})\,
\tfrac{\operatorname{vol}(B_d)}{(2\pi)^d}V\Lambda^d\,(1+o(1))]$ under a
cutoff.  The obstruction is the incompatibility of the identity embedding
with the Dyson-transported involution, not physical Hilbert-space
inequivalence.  The divergence of $r(k)$ itself is
a statement about the observable-space (phase-space) metric, which we
separate cleanly (Corollary~\ref{cor:scale}): in the canonically
normalized mode bundle its eigenvalues
$e^{\pm r(k)}\asymp(1+k^2)^{\pm1}$ give the two spectral directions
relative weights of orders $+2$ and $-2$ --- a Sobolev-pair structure
relative to that trivialization --- while the state-level mismatch
remains $O(1)$ per mode.

\paragraph{Verification.}
As in \cite{companion}, every finite-dimensional identity in this paper ---
the invariant formulas, the closed-form log-eigenvalues (including the
absolute-value form of \eqref{eq:xformula}), the PT ground-state
Gaussian and its Dyson-pullback identity \eqref{eq:pullback}, the
eigenvalue equation, the fidelity and occupation formulas, and
the limits $\sqrt3/2$ and $1/3$ --- is machine-verified symbolically, with
independent numerical checks (global optimization for
Theorem~\ref{thm:mindist}; high-precision quadrature for
Theorem~\ref{thm:groundstate}).  The artifact is the root of the
\texttt{area9innovation/weyl-gravity} repository of \cite{companion}
(scripts \texttt{symbolic/verify\_variational\_fock.py},
\texttt{symbolic/verify\_paper2\_referee.py},
\texttt{numeric/distortion\_scan.py}, \texttt{numeric/cartan\_checks.py};
SymPy 1.14.0, Python 3.14; cite the repository commit hash to fix the
artifact).

\section{The minimum-distortion theorem}\label{sec:mindist}

Work in the canonically normalized coordinates of \cite{companion}; all
adjoints refer to them.  The stabilizer is parametrized by two complex
angles,
\begin{equation}
C(\theta_1,\theta_2)=\sum_{j=1,2}\big(\cos\theta_j\,P_j+\sin\theta_j\,X_j\big),
\qquad X_j^2=-P_j ,
\end{equation}
with $P_j$ the projections onto the modes of $G_0'$ and $X_j$ the per-mode
complex structures; every admissible diagonalizer is $S=S_+C$.

\begin{definition}
For an admissible $S$ let
$\mathcal F(S)=\|\log(S^{\dg}S)\|_F^2$, the squared affine-invariant distance
of the observable-space metric $S^{\dg}S$ from the identity
\cite{Bhatia2007}.  For a stabilizer element $C=C_1\oplus C_2$ (blocks on the
two modes) define its \emph{Gram rapidities} $\tau_j\ge0$ by
\begin{equation}
W_j=C_jC_j^{\dg},\qquad \det W_j=1,\qquad
\cosh\tau_j=\tfrac12\tr W_j ,
\end{equation}
and set $a=\tfrac12(\tau_1+\tau_2)$, $b=\tfrac12(\tau_1-\tau_2)$.
\end{definition}

Note $\tau_1=\tau_2=0$ iff $W_1=W_2=I$ iff $C$ is unitary.  Since
$P=S^{\dg}S=C^{\dg}S_+^2C$ is Hermitian positive \emph{and} symplectic, its
eigenvalues form reciprocal pairs $e^{\pm x_1},e^{\pm x_2}$ with
$x_i\ge0$, and $\mathcal F(S)=2(x_1^2+x_2^2)$.

\begin{lemma}[Invariants]\label{lem:invariants}
Write $W=CC^{\dg}=W_1\oplus W_2$, $c=\cosh r$, $s=\sinh r$.  Decomposing
$W_j=\cosh\tau_j\,I+\sinh\tau_j\,(n_j\!\cdot\!\sigma)$ with real unit vectors
$n_j$, and letting $\chi\in[-1,1]$ denote the alignment
$n_1\!\cdot\!\tilde n_2$ (where $\tilde n_2$ is $n_2$ with its $\sigma_2$
component reflected),
\begin{align}
\cosh x_1+\cosh x_2 &= c\,(\cosh\tau_1+\cosh\tau_2), \label{eq:T}\\
\cosh x_1\cosh x_2 &= \tfrac12\Big[(2c^2-s^2)\cosh\tau_1\cosh\tau_2
+s^2\big(1+\sinh\tau_1\sinh\tau_2\,\chi\big)\Big]. \label{eq:Pi}
\end{align}
\end{lemma}

\begin{proof}
$\tr P=\tr(e^{rB}W)=c\,\tr W+s\,\tr(BW)$, and $\tr(BW)=0$ because $B$ is
off-diagonal and $W$ diagonal with respect to the mode splitting; this gives
\eqref{eq:T}.  For \eqref{eq:Pi}, expand
$\tr P^2=c^2\tr W^2+s^2\tr\big((BW)^2\big)$ (the cross terms again vanish
blockwise); the off-diagonal blocks of $B$ equal $-\sigma_2$, and the $2\times2$
identity $\sigma_2M\sigma_2=(\det M)\,M^{-\top}$ reduces
$\tr\big((BW)^2\big)=2\tr(\overline{W_2}^{\,-1}W_1)
=4\big[\cosh\tau_1\cosh\tau_2-\sinh\tau_1\sinh\tau_2\,\chi\big]$.
Converting $(\tr P,\tr P^2)$ to the symmetric functions of
$\cosh x_i$ yields \eqref{eq:Pi}.
\end{proof}

\noindent
(As verification evidence, both matrix identities are also checked against
random stabilizer elements in the repository.)

\begin{lemma}[Worst alignment; closed form]\label{lem:closedform}
At $\chi=-1$ the system \eqref{eq:T}--\eqref{eq:Pi} is solved exactly by
the unordered nonnegative pair
\begin{equation}\label{eq:xformula}
\{x_1,x_2\}=\big\{\varphi(b)+a,\ \big|\varphi(b)-a\big|\big\},
\qquad
\varphi(b)=\arccosh\big(\cosh r\,\cosh b\big)\ \ge\ r
\end{equation}
(the absolute value is required: for $a>\varphi(b)$ the signed expression
$\varphi-a$ is negative, while the $x_i$ are defined as nonnegative
log-eigenvalue exponents; equivalently
$\cosh x_{1,2}=\cosh(\varphi\pm a)$ with the principal nonnegative
$\arccosh$).
\end{lemma}

\begin{proof}
With $\alpha=\cosh a$, $\beta=\cosh b$ one finds from
\eqref{eq:T}--\eqref{eq:Pi} at $\chi=-1$, using
$\cosh\tau_1\cosh\tau_2\pm\sinh\tau_1\sinh\tau_2=\cosh(\tau_1\pm\tau_2)$,
\[
\cosh x_1+\cosh x_2=2c\,\alpha\beta,
\qquad
\cosh x_1\cosh x_2=\alpha^2+c^2\beta^2-1 ,
\]
whence
$\cosh x_{1,2}=c\alpha\beta\pm\sqrt{(\alpha^2-1)(c^2\beta^2-1)}
=\cosh(\varphi\pm a)$ with $\cosh\varphi=c\beta$, by the hyperbolic
addition formulas (elementary identities; also machine-checked in the
repository as verification evidence); since $\cosh$ is even, the
nonnegative exponents are $\varphi+a$ and $|\varphi-a|$.  Finally
$\cosh\varphi(b)-\cosh r=\cosh r(\cosh b-1)\ge0$ gives $\varphi(b)\ge r$ with
equality iff $b=0$.
\end{proof}

\begin{lemma}[Alignment monotonicity]\label{lem:monotone}
At fixed $\cosh x_1+\cosh x_2$, the quantity $x_1^2+x_2^2$ is nondecreasing
in $\cosh x_1\cosh x_2$; and \eqref{eq:Pi} is nondecreasing in $\chi$.
\end{lemma}

\begin{proof}
Parametrize $\cosh x_{1,2}=T/2\pm w$, $w\ge0$; the product decreases in $w$,
and
$\frac{d}{dw}(x_1^2+x_2^2)=2\big(\frac{x_1}{\sinh x_1}-\frac{x_2}{\sinh x_2}\big)
\le0$
because $x/\sinh x$ is strictly decreasing and $x_1\ge x_2$.  The second
statement is the sign of the coefficient
$s^2\sinh\tau_1\sinh\tau_2/2\ge0$ in \eqref{eq:Pi}.
\end{proof}

\begin{theorem}[Minimum distortion]\label{thm:mindist}
For every admissible diagonalizer $S=S_+C$,
\begin{equation}
\mathcal F(S)\ \ge\
4\Big[a^2+\arccosh^2\big(\cosh r\cosh b\big)\Big]\ \ge\ 4r^2=\mathcal F(S_+),
\end{equation}
with overall equality if and only if $C$ is unitary.  All minimizers share
the polar factor $S_+$, hence the same Hilbert-space metric
$\eta=e^{-Q}$: the Bender--Mannheim metric is the unique metric of any
minimum-distortion diagonalizer,
\begin{equation}
S_+\in\operatorname*{arg\,min}_{S^{\top}JS=J,\ S^{\top}G'S=G_0'}
d\big(I,S^{\dg}S\big),
\qquad
Q=-2\log\big(\text{polar factor of any minimizer}\big)\big|_{\text{metaplectic}} .
\end{equation}
\end{theorem}

\begin{proof}
By Lemma~\ref{lem:monotone}, lowering $\chi$ to $-1$ at fixed Gram rapidities
does not increase $x_1^2+x_2^2$; by Lemma~\ref{lem:closedform} the value there
is $\mathcal F=2\big[(\varphi+a)^2+(\varphi-a)^2\big]=4(\varphi^2+a^2)$;
and $\varphi(b)\ge r$.  Equality in the chain
forces $a=0$ and $b=0$, i.e.\ $\tau_1=\tau_2=0$, i.e.\ $C$ unitary; the
converse is immediate since then $S^{\dg}S$ is unitarily conjugate to
$S_+^2$.  If $C$ is unitary, $S_+C=S_+\cdot C$ is already a polar
decomposition, so the polar factor is $S_+$ and the metaplectic implementer
factorizes as $\rho'=U_C\,e^{-Q/2}$ with $U_C$ unitary, giving
$\eta'=\rho'^{\dg}\rho'=e^{-Q}$ for every minimizer.
\end{proof}

\begin{remark}
(i) The individual log-eigenvalues do \emph{not} satisfy $x_i\ge r$: numerical
scans show $\min_i x_i<r$ on open regions of the stabilizer; only the joint
quantity is bounded below, which is why the two-invariant argument is needed.
(ii) The equality set is the unitary part of the stabilizer, which in the
canonical coordinates is generically the finite group $\Z_2\times\Z_2$
\cite{companion}; abstractly it is contained in the maximal compact
$U(1)^2$ of $SO(2,\C)^2$.
(iii) The theorem is an instance of an orthogonal Cartan-norm projection
principle, proved in Section~\ref{ssec:cartan} together with its converse;
the elementary proof above remains valuable because it yields the exact
value \eqref{eq:xformula}, which the geometric argument does not.
\end{remark}

\subsection{Geometric recognition: an orthogonal Cartan-norm projection
principle}\label{ssec:cartan}

\paragraph{Normalization dictionary.}
Define the Cartan projection $\mu(S)$ of $S\in\Sp(4,\C)$ as the
chamber-ordered log-spectrum of $S^{\dg}S$ (twice the usual
log-singular-value vector), equip the Hermitian directions $\mathfrak p$
with the trace form $\langle\xi,\eta\rangle=\tr(\xi\eta)$, and normalize the
symmetric-space distance by $d_X(o,e^{\xi}\!\cdot o)=\|\xi\|_F$.  If
$\operatorname{spec}(\xi)=\{\pm a_1,\pm a_2\}$ then
$\mu(e^{\xi})=(2a_1,2a_2)$ (chamber-ordered), $\|\xi\|_F^2=2\sum_ia_i^2$,
and the three minimands used below are related by
\begin{equation}\label{eq:dictionary}
\boxed{\ \mathcal F(g)=2\|\mu(g)\|^2,
\qquad
\mathcal F(e^{\xi})=4\,d_X(o,e^{\xi}\!\cdot o)^2=\|2\xi\|_F^2 .\ }
\end{equation}
In particular $\min\sqrt{\mathcal F}=\|2\xi\|_F$,
$\min\|\mu\|=\|2\xi\|_F/\sqrt2$, and the first-variation coefficient $8$ in
\eqref{eq:firstvar} below is consistent with
$\mathcal F(e^{\xi})=4\|\xi\|_F^2$.  Throughout, $K$ denotes the
maximal compact subgroup determined by the physical inner product:
$K=\Sp(4,\C)\cap U(4)\cong\Sp(2)$ (the compact symplectic group, written
$USp(4)$ in part of the literature), and in the $n$-mode setting
$K=\Sp(2n,\C)\cap U(2n)\cong\Sp(n)$.  With these conventions
$\mu(S_+)=(r,r)$, which lies \emph{on the chamber wall} $x_1=x_2$ of type
$C_2$.

Theorem~\ref{thm:mindist} is thus a minimization of the Cartan-projection
\emph{norm} along a stabilizer coset.  We emphasize at the outset what is
classical and what is not.  The scalar projection inequality below is
standard symmetric-space geometry in the tradition of Mostow's self-adjoint
groups \cite{Mostow1955,BridsonHaefliger}: the compatible-subgroup orbit is
totally geodesic and convex, and normal geodesics realize nearest-point
projection.  (Kempf--Ness theory \cite{KempfNess1979} offers a related
variational perspective on minimal-norm points of complexified orbits, but
it is not the direct source of this coset inequality.)  The inequality
controls $\|\mu\|$, not the vector-valued Cartan projection in Weyl-chamber
order, and we do not claim it as new.  The structure specific to this
problem is: (i) the physically given stabilizer $H=\Stab(J,G_0')$ is
\emph{not} $\theta$-stable, so the classical theory does not see it
directly; (ii) its canonical $\theta$-stable hull is exactly the
mode-preserving $SL(2,\C)\times SL(2,\C)$
(Proposition~\ref{prop:hull}), through which minimality is inherited by
restriction; and (iii) the exact $\mathbb H^2\times\R$ value formula
\eqref{eq:xformula}.

Let $\mathbb P$ denote the cone of positive-definite Hermitian matrices with
the affine-invariant metric, a nonpositively curved (NPC) space in which
$t\mapsto d(\gamma(t),x)^2$ is strictly convex along nonconstant geodesics,
with $\frac{d}{dt}\,d(e^{t\zeta},P)^2\big|_{t=0}=-2\langle\zeta,\log
P\rangle$ \cite{Bhatia2007,BridsonHaefliger}.

\begin{definition}[Compatible subgroup]\label{def:compatible}
A closed subgroup $\mathcal C\subset GL(n,\C)$ is \emph{compatible} if it
is closed under conjugate transpose and admits the global Cartan
decomposition
\begin{equation}
\mathcal C=K_{\mathcal C}\exp\mathfrak c_{\mathfrak p},
\qquad
K_{\mathcal C}=\mathcal C\cap U(n),
\qquad
\mathfrak c_{\mathfrak p}
=\{\eta=\eta^{\dg}:e^{t\eta}\in\mathcal C\ \forall t\in\R\}.
\end{equation}
Every connected closed $\dg$-stable reductive matrix group is compatible;
mere closedness under $\dg$ is \emph{not} sufficient: the discrete group
$\mathcal C=\{A^n:n\in\Z\}$, $A=\operatorname{diag}(2,\tfrac12)
\in\Sp(2,\R)$, is closed and self-adjoint with all elements positive, yet
$A^{1/2}\notin\mathcal C$ and $\mathfrak c_{\mathfrak p}=0$, and it
violates the conclusions of Lemma~\ref{lem:totgeo} and
Theorem~\ref{thm:projection} below (take $\xi=-\tfrac34\log A$: then $c=A$
beats $c=e$).  Compatibility excludes such counterexamples.
\end{definition}

\begin{lemma}\label{lem:totgeo}
Let $\mathcal C\subset GL(n,\C)$ be a compatible subgroup and let
$\mathbb P_{\mathcal C}=\mathcal C\cap\mathbb P$ be its set of
positive elements.  Then
$\mathbb P_{\mathcal C}=\exp\mathfrak c_{\mathfrak p}$ is a topologically
closed, geodesically closed (hence convex) subset of $\mathbb P$ --- so it
admits unique metric projection from any point of $\mathbb P$ --- and
$\mathbb P_{\mathcal C}=\{cc^{\dg}:c\in\mathcal C\}$.
\end{lemma}

\begin{proof}
A positive element of $\mathcal C$ has, by the global decomposition and
uniqueness of the polar decomposition, trivial unitary part, so
$\mathbb P_{\mathcal C}=\exp\mathfrak c_{\mathfrak p}$; in particular
$A\in\mathbb P_{\mathcal C}$ has $\log A\in\mathfrak c_{\mathfrak p}$ and
$A^{t}=e^{t\log A}\in\mathcal C$ for all real $t$.  For
$A,B\in\mathbb P_{\mathcal C}$ the geodesic is
$t\mapsto A^{1/2}\big(A^{-1/2}BA^{-1/2}\big)^{t}A^{1/2}$
\cite{Bhatia2007}; it lies in $\mathcal C$ because $A^{\pm1/2}$, the
positive matrix $M=A^{-1/2}BA^{-1/2}\in\mathbb P_{\mathcal C}$, and its
real powers $M^{t}$ all do, and it consists of positive elements, hence
stays in $\mathbb P_{\mathcal C}$.  For the Gram-set claim:
$cc^{\dg}=k\,e^{2\eta}k^{-1}=e^{2\,\mathrm{Ad}_k\eta}
\in\exp\mathfrak c_{\mathfrak p}$ for $c=k e^{\eta}$
($\mathrm{Ad}_{K_{\mathcal C}}$ preserves $\mathfrak c_{\mathfrak p}$),
and conversely $A=(A^{1/2})(A^{1/2})^{\dg}$ with
$A^{1/2}\in\mathcal C$.  Topological closedness of
$\mathbb P_{\mathcal C}=\mathcal C\cap\mathbb P$ is immediate from
closedness of $\mathcal C$; unique metric projection onto closed convex
subsets is standard in NPC spaces \cite{BridsonHaefliger}.
\end{proof}

\begin{theorem}[Orthogonal Cartan-norm projection principle]
\label{thm:projection}
Let $\mathcal C\subset GL(n,\C)$ be a compatible subgroup
(Definition~\ref{def:compatible}) and let $\xi=\xi^{\dg}$.  Then:
\begin{enumerate}
\item[(a)] For every $\eta\in\mathfrak c_{\mathfrak p}$,
\begin{equation}\label{eq:firstvar}
\frac{d}{dt}\Big|_{t=0}\,
\big\|\log\big((e^{\xi}e^{t\eta})^{\dg}e^{\xi}e^{t\eta}\big)\big\|_F^2
=8\,\langle\eta,\xi\rangle .
\end{equation}
\item[(b)] If $\xi\perp\mathfrak c_{\mathfrak p}$, then for every
$c\in\mathcal C$,
\begin{equation}
\big\|\log\big((e^{\xi}c)^{\dg}e^{\xi}c\big)\big\|_F\ \ge\ \|2\xi\|_F,
\qquad
\operatorname*{arg\,min}_{c\in\mathcal C}
\big\|\log\big((e^{\xi}c)^{\dg}e^{\xi}c\big)\big\|_F
=\mathcal C\cap U(n) .
\end{equation}
\item[(c)] Conversely, if the tangential part
$\xi_T=\mathrm{proj}_{\mathfrak c_{\mathfrak p}}\,\xi$ is nonzero, then
$c=e$ is not a minimizer: by \eqref{eq:firstvar} with $\eta=-\xi_T$ the
functional decreases at rate $8\|\xi_T\|^2$.  Hence
\begin{equation}
e\ \text{minimizes}\ c\mapsto
\big\|\log\big((e^{\xi}c)^{\dg}e^{\xi}c\big)\big\|_F
\ \text{on}\ \mathcal C
\quad\Longleftrightarrow\quad
\xi\perp\mathfrak c_{\mathfrak p} .
\end{equation}
\end{enumerate}
\end{theorem}

\begin{proof}
By congruence invariance of the affine metric,
$d\big(I,(e^{\xi}c)^{\dg}e^{\xi}c\big)=d\big((cc^{\dg})^{-1},e^{2\xi}\big)$,
and $(cc^{\dg})^{-1}=dd^{\dg}$ with $d=c^{-\dg}\in\mathcal C$, so it
ranges over $\mathbb P_{\mathcal C}$ as $c$ ranges over $\mathcal C$
(Lemma~\ref{lem:totgeo}).  For (a), take $c=e^{t\eta}$: then
$(cc^{\dg})^{-1}=e^{-2t\eta}$ is a geodesic through $I$ with velocity
$-2\eta$, and the NPC first-variation formula gives
$\frac{d}{dt}d(e^{-2t\eta},e^{2\xi})^2\big|_0
=-2\langle-2\eta,2\xi\rangle=8\langle\eta,\xi\rangle$.
For (b), every point of $\mathbb P_{\mathcal C}$ is joined to $I$ by a
geodesic $t\mapsto e^{2t\eta}$, $\eta\in\mathfrak c_{\mathfrak p}$, inside
$\mathbb P_{\mathcal C}$; along it $h(M)=d(M,e^{2\xi})^2$ is strictly convex
unless constant, and its first variation at $I$ vanishes by orthogonality.
Hence $I$ is the unique minimizer of $h$ on $\mathbb P_{\mathcal C}$, with
value $\|2\xi\|_F^2$; the minimizing $c$ are exactly those with
$cc^{\dg}=I$.  Part (c) is immediate from (a).
\end{proof}

The numerical experiments that preceded this proof (random tangential and
normal directions $\xi$) are retained in the repository as regression
checks of \eqref{eq:firstvar} and of both directions of (c).

\begin{proposition}[Canonical compatible hull, stratified]\label{prop:hull}
Let $\theta(g)=(g^{\dg})^{-1}$ and let
$C_\theta(H)$ be the smallest closed subgroup containing
$H=\Stab(J,G_0')$ and $\theta(H)$.  In the canonical coordinates,
\begin{equation}
C_\theta(H)=C_1\times C_2,
\qquad
C_j=\begin{cases}
SL(2,\C), & \lambda_j\neq1,\\[1mm]
SO(2,\C), & \lambda_j=1 ,
\end{cases}
\end{equation}
acting on the $j$-th mode plane; here $\lambda_j>0$ is the dimensionless
mode-aspect parameter of the canonical coordinates
(\cite{companion}; $\lambda_1=\w_2$, $\lambda_2=\w_1$ in normalized units
for the symmetric split, with the split-invariant
$\lambda_1/\lambda_2=\w_2/\w_1<1$).  On the generic stratum the compatible
hull is the full mode-preserving symplectic subgroup: not merely a
convenient supergroup but the minimal $\theta$-stable subgroup forced by
the physical stabilizer.  At an alignment value $\lambda_j=1$ the $j$-th
factor of $H$ is already $\theta$-stable and the hull degenerates to $H_j$
itself.
\end{proposition}

\begin{proof}
$H$ and $\theta(H)$ are contained in the block-diagonal subgroup, so
$C_\theta(H)\subseteq SL(2,\C)^2$.  For the converse it suffices to show the
real Lie algebra generated by
$\mathfrak h_j=\C X_j$ and $\theta(\mathfrak h_j)=\C X_j^{\dg}$ is all of
$\mathfrak{sl}(2,\C)$ on each mode.  On mode $1$ (block coordinates),
$X_1=\big(\begin{smallmatrix}0&\lambda_1\\-\lambda_1^{-1}&0\end{smallmatrix}\big)$
(with $\lambda_1=\w_2$ in normalized units for the symmetric split) and
\begin{equation}
[X_1,X_1^{\dg}]=(\lambda_1^2-\lambda_1^{-2})
\begin{pmatrix}1&0\\0&-1\end{pmatrix}\neq0
\qquad(\lambda_1\neq1),
\end{equation}
and the real span of
$X_1,X_1^{\dg},[X_1,X_1^{\dg}]$ and their $i$-multiples is six-dimensional,
hence equals $\mathfrak{sl}(2,\C)$; likewise on mode $2$ with $\lambda_2$.
The generated group is then the connected $SL(2,\C)$ on each generic mode.
At $\lambda_j=1$ one has $X_j^{\dg}=-X_j$, so
$\theta(\mathfrak h_j)=\mathfrak h_j$ and the closure is
$\mathfrak{so}(2,\C)$ itself.
\end{proof}

\noindent
(As regression evidence, not as part of the proof: the Lie-closure
dimensions $6$ and $2$ are also confirmed computationally in the
repository.)

\begin{remark}[Alignment: hull contraction versus physical enhancement]
\label{rem:enhancement}
The stratification of Proposition~\ref{prop:hull} is mirrored by two
distinct discontinuities at the alignment locus $\lambda_j=1$, which must
not be conflated.  For the
\emph{hull}, the compact locus contracts, $C_j\cap K:\ SU(2)\to U(1)$,
because the hull itself shrinks from $SL(2,\C)$ to $SO(2,\C)$; accordingly,
for minimization over the hull the equality locus \emph{changes} from
$SU(2)$-type to $U(1)$-type.  For the \emph{physical stabilizer} the
unitary locus is enhanced, $H_j\cap K:\ \Z_2\to U(1)$, and it is only for
the variational problem restricted to $H$ that the equality locus of
Theorem~\ref{thm:mindist} \emph{enlarges} at alignment.  The enhancement is
a property of the physical embedding (the twist), not of the abstract
groups, and marks the parameter values at which the twisted and compatible
pictures coincide.
\end{remark}

\begin{corollary}[Second proof of Theorem~\ref{thm:mindist}]
\label{cor:secondproof}
Apply Theorem~\ref{thm:projection}(b) with $\mathcal C$ the stratified
compatible hull $C_\theta(H)$ of Proposition~\ref{prop:hull}: on the
generic stratum $\mathcal C=SL(2,\C)^2$ --- connected, $\dg$-stable and
reductive, hence compatible --- and at alignment one or both factors
reduce to $SO(2,\C)$; orthogonality of $\xi$ holds on every stratum
because $\xi$ is block-off-diagonal.  On the generic stratum the positive
part $\mathbb P_{\mathcal C}\cong\mathbb H^3\times\mathbb H^3$ is totally
geodesic in $\mathbb P$, the Gram data $(\tau_j,n_j)$ of
Lemma~\ref{lem:invariants} being polar coordinates on the two hyperbolic
factors; take $\xi=\tfrac r2B$, which is block-off-diagonal and hence
orthogonal to $\mathfrak c_{\mathfrak p}$ (the per-mode traceless Hermitian
blocks).  This yields $\mathcal F(S_+c)\ge\mathcal F(S_+)$ for every $c$ in the
hull, with the two equality records kept separate:
\begin{equation}
\operatorname*{arg\,min}_{c\in\mathcal C}\mathcal F(S_+c)=\mathcal C\cap K,
\qquad
\operatorname*{arg\,min}_{h\in H}\mathcal F(S_+h)=H\cap K ,
\end{equation}
the second following from the first because $e\in H\subset\mathcal C$
realizes the unrestricted minimum.
\end{corollary}

\begin{remark}\label{rem:pythagoras}
(i) The closed form \eqref{eq:xformula} is the Pythagorean law of the
ambient symmetric space: the worst-alignment configurations sweep a totally
geodesic $\mathbb H^2\times\R$, on which
$\arccosh(\cosh r\cosh b)$ is the hyperbolic hypotenuse law
($\cosh h=\cosh r\cosh b$, legs meeting at a right angle in $\mathbb H^2$)
while $a$ adds by flat Pythagoras in the $\R$ factor.  This exact value is
content beyond the projection principle, which yields only the inequality.
(ii) On the compactness subtlety: $H\cap K\cong\Z_2\times\Z_2$ refers to the
intersection with the ambient maximal compact determined by the physical
inner product.  Abstractly $SO(2,\C)\cong\C^{\times}$ has maximal compact
$U(1)$ after a suitable conjugation; the finiteness is a property of the
physical embedding --- which is precisely the point: the embedding, not the
group, carries the obstruction, and Proposition~\ref{prop:hull} locates the
minimal compatible geometry it forces.
\end{remark}

\subsection{The $n$-mode statement}\label{ssec:nmode}

The projection principle is dimension-independent, so the variational
selection extends verbatim to any number of PT mode pairs.

\begin{theorem}[$\Sp(2n,\C)$ selection principle]\label{thm:nmode}
Fix a symplectic splitting of $\C^{2n}$ into $n$ canonical mode planes and
let $\mathcal C_n\cong\prod_{j=1}^{n}SL(2,\C)\subset\Sp(2n,\C)$ be the
mode-preserving subgroup (connected and $\dg$-stable, hence compatible),
with positive part $\mathbb P_{\mathcal C_n}\cong(\mathbb H^3)^{n}$.  The orthogonal complement
of $\mathfrak c_{n,\mathfrak p}$ in the Hermitian Hamiltonian directions
consists exactly of the \emph{inter-mode} couplings (Hermitian Hamiltonian
matrices with vanishing mode-diagonal blocks).  Consequently, for every
inter-mode $\xi=\xi^{\dg}$,
\begin{equation}
\min_{c\in\mathcal C_n}
\big\|\log\big((e^{\xi}c)^{\dg}e^{\xi}c\big)\big\|_F=\|2\xi\|_F,
\qquad
\operatorname*{arg\,min}=\mathcal C_n\cap K,
\end{equation}
and for any quadratic PT Hamiltonian whose positive normal-form
diagonalizer is $e^{\xi}$ with inter-mode $\xi$, and whose stabilizer is
contained in $\mathcal C_n$, the positive diagonalizer is the unique
minimum-distortion diagonalizer up to matrix-unitary stabilizer elements.
\end{theorem}

\begin{proof}
Write $\xi\in\mathfrak p$ in $2\times2$ blocks $(\xi_{ij})_{i,j=1}^{n}$
relative to $V=V_1\oplus\dots\oplus V_n$.  A one-parameter group
$e^{t\eta}$ lies in $\mathcal C_n$ iff $\eta$ is block-diagonal, so
$\mathfrak c_{n,\mathfrak p}=\{\xi\in\mathfrak p:\xi_{ij}=0,\ i\neq j\}$
(per-mode Hermitian traceless blocks: Hermiticity and the $2\times2$
Hamiltonian condition make each $\xi_{jj}$ a traceless Hermitian block).
Since $\tr(\xi\eta)=\sum_{ij}\tr(\xi_{ij}\eta_{ji})$ and $\eta$ is
block-diagonal, the pairing sees only $\sum_j\tr(\xi_{jj}\eta_{jj})$, so
$\xi\perp\mathfrak c_{n,\mathfrak p}$ iff every diagonal block $\xi_{jj}$
vanishes: the normal space is exactly the inter-mode couplings, whose
paired off-diagonal blocks $\xi_{ji}$ are determined from $\xi_{ij}$ by the
Hermitian and symplectic conditions.  The dimension identity
$n(2n+1)=3n+2n(n-1)$ is then a check, not a step.  Everything else is
Theorem~\ref{thm:projection} applied to $\mathcal C_n$.
\end{proof}

The theorem ends deliberately at the finite-dimensional Gram-matrix
statement.  Matrix unitarity is not automatically quantum unitarity: a
unitary operator implementing $UVU^{-1}=CV$ in quadrature coordinates
($V=V^{\dg}$) must send self-adjoint canonical variables to self-adjoint
canonical variables, which requires $C$ to preserve the physical real
structure, not merely $C^{\dg}C=I$ --- a generic element of
$\Sp(2n,\C)\cap U(2n)$ has complex coefficients and defines no
$*$-automorphism of the real CCR algebra.  The Hilbert-space conclusion
therefore needs an explicit hypothesis:

\begin{corollary}[Hilbert-metric constancy under $*$-compatibility]
\label{cor:starcompat}
In the setting of Theorem~\ref{thm:nmode}, assume in addition that every
minimizer $C\in\mathcal C_n\cap K$ arising from the physical stabilizer
preserves the selected CCR real structure (i.e.\ is real in quadrature
coordinates, hence metaplectically implemented by a genuine unitary).
Then all minimizers define the same Hilbert-space metric, by the
polar-factor argument of Theorem~\ref{thm:mindist}.  For the PU
oscillator this hypothesis holds: the unitary stabilizer elements in the
canonical coordinates are the real sign matrices
$\{\pm P_1\pm P_2\}\cong\Z_2^2$ generically, and ordinary real rotations
at alignment \cite{companion}, so the PU metric-constancy conclusion of
Theorem~\ref{thm:mindist} is unaffected.  For a general $n$-mode system
the hypothesis is to be verified per system, alongside the containment
of the stabilizer in $\mathcal C_n$.
\end{corollary}

For the PU oscillator, $n=2$ and $\xi=\tfrac r2B$ is inter-mode; the
Bender--Mannheim metric is the case in point.  The model-dependent input in
any application is thus reduced to a single check: that the positive
diagonalizer direction is inter-mode for the chosen splitting, i.e.,
trace-orthogonal to $\mathfrak c_{n,\mathfrak p}$ --- and by
Theorem~\ref{thm:projection}(c) this check is not only sufficient but
necessary.  We emphasize that the canonical-hull computation of
Proposition~\ref{prop:hull} is proved here for the two-mode PU stabilizer;
for a general $n$-mode system the containment of the physical stabilizer in
$\mathcal C_n$, and the structure of its compatible hull across mixed
exceptional strata, are hypotheses to be verified per system.

\section{The free fourth-order field}\label{sec:field}

\subsection{Modewise data and the phase-space Hilbert scale}

For $(\Box+m_1^2)(\Box+m_2^2)\phi=0$, each spatial momentum $k$ carries a PU
pair with $\w_i(k)=\sqrt{k^2+m_i^2}$ and $m_1>m_2>0$.  Since
$(\w_1-\w_2)(\w_1+\w_2)=m_1^2-m_2^2$ identically, the rapidity obeys
\eqref{eq:rk} \emph{exactly}, interpolating between
$r(0)=\log\frac{m_1+m_2}{m_1-m_2}$ and $r(k)=2\log|k|+O(1)$.

\begin{corollary}[Phase-space Hilbert scale]\label{cor:scale}
In the canonically normalized modewise coordinates, the observable-space
metric $M_{\mathrm{obs}}(k)=S_+(k)^{\dg}S_+(k)$ has eigenvalues
$e^{\pm r(k)}$ with
\begin{equation}
e^{r(k)}=\frac{\big(\w_1(k)+\w_2(k)\big)^2}{m_1^2-m_2^2}
\asymp 1+k^2 ,
\qquad
e^{-r(k)}\asymp(1+k^2)^{-1} .
\end{equation}
On the spectral subspaces $\Pi_\pm(k)$ of $B$, \emph{in the canonically
normalized mode bundle}, the two spectral directions therefore carry
relative weights of orders $+2$ and $-2$: the induced norm is not
uniformly equivalent to the standard one, and no scalar renormalization
makes it so.  We emphasize that this is a statement relative to the
$k$-dependent normalized trivialization
($d_x(k)d_y(k)=\gamma\w_1(k)\w_2(k)$); an invariant Sobolev
classification of the \emph{original} field variables would require the
pullback $D(k)^{\dg}M_{\mathrm{obs}}(k)D(k)$ in a fixed Fourier-space
trivialization, since part of the weight may be carried by the
normalization itself.  We use only the normalized-bundle statement below.
\end{corollary}

This is the precise content of the observation that the modewise metric
``blows up in the UV''.  It concerns the classical (phase-space) metric.
The quantum state-level comparison is finer, and is computed next.

\subsection{The PT ground state, exactly}

Fix a mode and work in the normalized coordinates $(x,y)$ with
$\w_1>\w_2>0$.  The positive Hamiltonian $\widetilde H'$ has ground state
$\phi_0\propto e^{-\frac12(x^2/\w_2+\w_1y^2)}$.  The PT Hamiltonian in these
coordinates is the differential operator
\begin{equation}
H_{\mathrm{PT}}'=\tfrac12\Big[\tfrac{\w_1^2+\w_2^2}{\w_1\w_2}\,x^2
+\w_1\w_2\,y^2-\w_1\w_2\,\partial_x^2\Big]-x\,\partial_y .
\end{equation}

\begin{theorem}[PT ground state]\label{thm:groundstate}
The unique Gaussian solution of the modewise annihilation equations
$\big(\ell_j^{\top}S_+^{-1}V\big)\psi_0=0$ is
\begin{equation}
\psi_0(x,y)\ \propto\
\exp\!\Big[-\tfrac12\Big(\tfrac{\w_1+\w_2}{\w_1\w_2}\,x^2
+2xy+(\w_1+\w_2)\,y^2\Big)\Big],
\end{equation}
a \emph{real} Gaussian with
$\det(\mathrm{quadratic\ form})=\frac{\w_1^2+\w_1\w_2+\w_2^2}{\w_1\w_2}>0$:
normalizable for all $\w_1>\w_2>0$, and
$H_{\mathrm{PT}}'\psi_0=\tfrac{\w_1+\w_2}{2}\,\psi_0$.  Moreover
$\psi_0$ is exactly the Dyson pullback of the normal-form vacuum,
\begin{equation}\label{eq:pullback}
\psi_0=\rho^{-1}\phi_0
\qquad(\text{up to normalization};\ \rho=e^{-Q/2}),
\end{equation}
since $\rho^{-1}V\rho=S_+^{-1}V$ transports the annihilation conditions
of $\phi_0$ into those defining $\psi_0$ (machine-verified at the
Gaussian level).  Its \emph{identity-embedding Gaussian fidelity} with
the two-oscillator vacuum --- both wavefunctions normalized in the same
auxiliary standard $L^2$ norm, canonical variables identified by the
identity map --- and its \emph{auxiliary standard-Fock occupation}
relative to it are exactly
\begin{align}
f(\w_1,\w_2)&=\big|\langle\hat\phi_0|\hat\psi_0\rangle\big|^2
=\frac{4\w_1\sqrt{\w_1^2+\w_1\w_2+\w_2^2}}{4\w_1^2+3\w_1\w_2+\w_2^2},\\[2pt]
\langle N\rangle(\w_1,\w_2)
&=\langle\hat\psi_0|(N_1+N_2)|\hat\psi_0\rangle
=\frac{\w_2(\w_1^2+\w_2^2)}{2\w_1(\w_1^2+\w_1\w_2+\w_2^2)},
\qquad \langle N_1\rangle=\langle N_2\rangle .
\end{align}
In particular, along the field dispersion $\w_i=\w_i(k)$,
\begin{equation}\label{eq:universal}
\lim_{|k|\to\infty}f=\frac{\sqrt3}{2}\approx0.866,
\qquad
\lim_{|k|\to\infty}\langle N\rangle=\frac13 ,
\end{equation}
universal constants independent of $m_1,m_2,\gamma$.
\end{theorem}

\begin{proof}
Solving the two first-order equations for a Gaussian ansatz gives the stated
quadratic form (one solution); the eigenvalue equation is a direct
substitution.  The fidelity is the standard two-Gaussian determinant formula;
the occupations follow from covariance algebra,
$\langle N_j\rangle=\tfrac12 c_j^{\top}\Sigma c_j$ with
$\Sigma$ the covariance of $|\psi_0|^2$ and $c_j$ the coefficient vectors of
$a_j\psi_0$.  The limits follow with $\w_1/\w_2\to1$.
(Independent symbolic verification of each step, and a 25-digit quadrature
check of the fidelity, are recorded in the repository as verification
evidence.)
\end{proof}

\begin{remark}
Two features deserve emphasis.  (i) $\psi_0$ is a smooth, normalizable
Gaussian even as $\w_1\to\w_2$: the exceptional-point pathology
(divergence of $S_+$, Jordan structure) lives in the similarity
transformation and the excited spectrum, not in the ground state itself.
(ii) A naive Bogoliubov estimate would assign the modewise map a squeezing
$\sim\sinh^2\!\big(r(k)/2\big)\sim k^2$ and hence a violently divergent
particle content per mode.  The exact result \eqref{eq:universal} is
$O(1)$: the diverging rapidity $r(k)$ is carried almost entirely by
directions that do not create quanta.  The naive estimate fails because
$\rho=e^{-Q/2}$ is not $*$-preserving on the original Fock structure ---
the transformed ``annihilators'' $\rho a_j\rho^{-1}$ do not satisfy the
unitary Bogoliubov normalization ($\alpha\alpha^{\dg}-\beta\beta^{\dg}
=\cosh r\ne1$) --- so Shale--Stinespring-type criteria cannot be applied to
$\rho$ directly.  Comparing the two \emph{states} as vectors of the same
auxiliary $L^2$ avoids this trap --- at the price, made explicit below,
that the comparison is then the auxiliary identity-embedding one, not the
physical $\eta$-completion.
\end{remark}

\subsection{The physical completion: pointed unitary equivalence}

\begin{proposition}[Pointed unitary equivalence of the physical
completions]\label{prop:pointed}
Equip the mode space with the physical inner product
$\langle u,v\rangle_{\eta_k}=\langle u,\eta_kv\rangle_{\mathrm{std}}$.
Then $U_k:=\rho_k$ is unitary from the PT mode space to the standard mode
space by construction, and by \eqref{eq:pullback} it is \emph{pointed}:
\begin{equation}
U_k\psi_k=\phi_k,
\qquad
\langle\phi_k,U_k\psi_k\rangle=1
\quad\text{for every mode }k .
\end{equation}
Consequently the pointed infinite tensor product
$U=\bigotimes_kU_k$ exists with no von Neumann obstruction and maps
$\bigotimes_k\psi_k\mapsto\bigotimes_k\phi_k$: \emph{with the
$\eta$-inner product and the Dyson-transported observable algebra
$A_{\mathrm{std}}=\rho A_{\mathrm{PT}}\rho^{-1}$, the physical PT
completion of the free fourth-order field is globally unitarily
equivalent to the positive normal-form Fock completion},
\begin{equation}
\omega_{\mathrm{PT}}(A_{\mathrm{PT}})
=\frac{\langle\psi,\eta A_{\mathrm{PT}}\psi\rangle_{\mathrm{std}}}
       {\langle\psi,\eta\psi\rangle_{\mathrm{std}}}
=\langle\phi,A_{\mathrm{std}}\phi\rangle_{\mathrm{std}} .
\end{equation}
In particular the physical expectation of the transported number
observable vanishes:
$\langle\rho^{-1}N\rho\rangle_{\eta}=\langle\phi,N\phi\rangle=0$.  The
unboundedness of $\rho$ relative to the standard norm
(Corollary~\ref{cor:scale}) is real but does not obstruct this
equivalence.
\end{proposition}

The $\sqrt3/2$ and $1/3$ of Theorem~\ref{thm:groundstate} are therefore
\emph{not} the physical overlap and particle content of the PT vacuum;
they quantify the mismatch of two auxiliary Gaussian complex structures
when the canonical variables are identified by the identity map.  That
mismatch is itself an exact and structured obstruction:

\subsection{The auxiliary identity-embedding obstruction}

\begin{theorem}[Auxiliary identity-embedding obstruction]
\label{thm:disjoint}
Let $m_1>m_2>0$, $d\ge1$, and consider the free fourth-order field in a
box of volume $V$ with momentum cutoff $\Lambda$.  Regard the normalized
wavefunctions $\hat\psi_0(k)$ and $\hat\phi_0(k)$ as vectors of one
auxiliary Schr\"odinger representation --- same standard $L^2$ inner
product, same standard CCR involution, canonical variables identified by
the identity map --- and let
$\Psi_{\mathrm{aux}}=\bigotimes_k\hat\psi_0(k)$,
$\Omega_{\mathrm{aux}}=\bigotimes_k\hat\phi_0(k)$ be the two auxiliary
product states.  Then (independent modes labelled by the $(k,-k)$ pairs
of the real field; a different mode convention changes the prefactor
below by a fixed factor):
\begin{enumerate}
\item the auxiliary overlap has the exact leading asymptotics
\begin{equation}
-\log\big|\langle\Omega_{\mathrm{aux}}|\Psi_{\mathrm{aux}}\rangle\big|^2
=\Big[\log\tfrac{2}{\sqrt3}\Big]\,
\frac{\operatorname{vol}(B_d)}{(2\pi)^d}\,V\Lambda^d
+o\big(V\Lambda^d\big),
\end{equation}
with universal leading coefficient (since $f(k)\to\sqrt3/2$ and
$-\tfrac1V\log|\langle\Omega,\Psi\rangle|^2
=-\int_{|k|\le\Lambda}\tfrac{d^dk}{(2\pi)^d}\log f(k)$);
\item the auxiliary standard-Fock occupation density diverges in the UV:
$\frac1V\langle N_{\mathrm{tot}}\rangle
=\int^{\Lambda}\!\frac{d^dk}{(2\pi)^d}\,\langle N\rangle(k)
\sim\frac13\int^{\Lambda}\!\frac{d^dk}{(2\pi)^d}\to\infty$ as
$\Lambda\to\infty$, in \emph{every} dimension $d\ge1$;
\item in the thermodynamic/continuum limit the incomplete-tensor-product
classes of $\Psi_{\mathrm{aux}}$ and $\Omega_{\mathrm{aux}}$ differ
($\sum_k(1-|\langle\hat\psi_0|\hat\phi_0\rangle|)=\infty$), so the
standard-$L^2$ Gaussian state defined by the PT eigenfunction is
disjoint from the standard two-field vacuum state \emph{under the
identity embedding of the auxiliary canonical variables}: no density
matrix in the auxiliary Fock representation reproduces it, and vice
versa.
\end{enumerate}
This theorem compares two vector states on one fixed auxiliary standard
CCR algebra.  It does \emph{not} compare the physical $\eta$-completion
with the transported normal-form completion --- those are equivalent by
Proposition~\ref{prop:pointed} --- and its conclusion may not be called
disjointness of the physical PT representation.
\end{theorem}

\begin{proof}
(1) and (2) are immediate from Theorem~\ref{thm:groundstate}: $f(k)$ and
$\langle N\rangle(k)$ are continuous in $k$ with the universal limits
\eqref{eq:universal}; the Riemann-sum limit of $-\log f$ gives the stated
coefficient.  (3) is von Neumann's
equivalence criterion for incomplete tensor products \cite{vN1939} applied
to the two product vectors, combined with purity: inequivalent pure product
states generate disjoint representations.
\end{proof}

\begin{remark}[What this does and does not say]\label{rem:meaning}
The two results together resolve what would otherwise be a contradiction.
The physical Hilbert space, defined as the completion of the
modewise-unitarized theory, is naturally identified with the Fock space
of the free pair $(m_1,m_2)$, globally and unitarily
(Proposition~\ref{prop:pointed}); on it the dynamics is unitary, exactly
as argued by Bender and Mannheim modewise --- and this equivalence is
\emph{pointed}, with per-mode overlap exactly one.  What fails is
something different: the identity embedding of the auxiliary canonical
variables does not extend to a normal equivalence of the two auxiliary
quasifree states (Theorem~\ref{thm:disjoint}), with an $O(1)$ discrepancy
per UV mode.  The obstruction is the incompatibility between the
\emph{identity embedding of canonical variables} and the
\emph{Dyson-transported physical involution} --- an involution-level
mismatch, not a Hilbert-space inequivalence, in line with the three-level
distinction (complex canonical dynamics $\neq$ real form/involution
$\neq$ Hilbert completion) developed for the field theory in the
companion papers.  Practical consequences survive in sharpened form: any
interacting construction must declare which embedding of the canonical
variables it uses; formal manipulations that insert $\eta=e^{-Q}$ as if
it were a bounded operator on the covariant Fock space are unsupported
(by Corollary~\ref{cor:scale} the metric is unbounded with unbounded
inverse); and PT-based unitarity claims for quadratic gravity
\cite{QQG2024} require an explicit statement of involution and
completion, not only of the modewise similarity.
\end{remark}

\section{Discussion}\label{sec:discussion}

\paragraph{The variational principle.}
Theorem~\ref{thm:mindist} upgrades the uniqueness theorem of
\cite{companion}: the Bender--Mannheim metric is not merely the unique
positive point of the diagonalizer coset but its unique distortion minimizer,
with the exact excess $4[a^2+\arccosh^2(\cosh r\cosh b)]-4r^2$ in the Gram
rapidities of the stabilizer.  Section~\ref{ssec:cartan} shows this is an
instance of an orthogonal Cartan-norm projection principle along
$\dg$-closed subgroups --- classical in spirit
\cite{Mostow1955,KempfNess1979}, here proved in if-and-only-if form with
the minimal compatible hull identified canonically
(Proposition~\ref{prop:hull}) --- while the elementary proof supplies the
minimizing value in closed form, which the abstract theory does not.  The
$n$-mode version (Theorem~\ref{thm:nmode}) makes the portability to
general quadratic PT systems a theorem rather than a claim.

A scope distinction now matters: minimum distortion is \emph{free
geometric optimality}, not \emph{interacting dynamical stability}.
Among admissible free positive diagonalizers the Bender--Mannheim metric
is the geometrically least distorting; the theorem does not make it the
most stable under interaction, uniquely physical, or guaranteed to
survive scattering.  Indeed the companion interaction paper
\cite{companionV} proves that the order-by-order deformation of this
metric is obstructed at rational resonances and on open scattering
shells.  That result sharpens rather than weakens the present one: the
obstruction strikes the \emph{optimal} free positive metric, so the
interacting failure cannot be dismissed as the consequence of an
arbitrary or badly chosen metric.

\paragraph{The three levels, again.}
The field-theoretic analysis separates three statements that are easily
conflated.  The \emph{phase-space} metric degenerates in the UV exactly as
the rapidity formula \eqref{eq:rk} dictates ($e^{\pm r(k)}\asymp
k^{\pm2}$ in the normalized mode bundle); this is a statement about
observables.  The \emph{physical completion} is globally unitarily
equivalent to the normal-form Fock completion, by a pointed unitary with
per-mode overlap one (Proposition~\ref{prop:pointed}); this is a
statement about Hilbert completions.  The \emph{auxiliary
identity-embedding} mismatch saturates at universal constants
($\sqrt3/2$, $1/3$) and its product states are disjoint
(Theorem~\ref{thm:disjoint}); this is a statement about the
incompatibility of the identity embedding with the Dyson-transported
involution.  The revised central conclusion is therefore:
\begin{quote}
\emph{The Dyson map gives an exact pointed unitary equivalence between
the physical PT completion and the positive normal-form Fock completion.
However, if the two Gaussian vacua are compared through the identity
embedding of the auxiliary canonical variables and the standard
involution, their standard-CCR product states are disjoint, with
universal ultraviolet fidelity and occupation defects.}
\end{quote}
This aligns with the three-level structure --- complex canonical dynamics
$\neq$ real form/involution $\neq$ Hilbert completion --- used in the
companion papers.  All quantities are exact consequences of the same
modewise data.  The constants are universal because the UV limit
$\w_1/\w_2\to1$ erases all mass and coupling dependence --- remarkably, the
same limit in which the finite-dimensional theory degenerates (the
exceptional point) is the limit in which the ground-state data become
featureless.

\paragraph{Outlook.}
Three natural continuations: (i) sharpen Theorem~\ref{thm:nmode} into a
classification: for which quadratic PT Hamiltonians is the positive
diagonalizer direction inter-mode for some canonical splitting
(by Theorem~\ref{thm:projection}(c) this is necessary and sufficient), and
what replaces the selection principle when it fails; (ii) determine
whether the auxiliary identity-embedding obstruction
survives interactions in low dimensions, where the free-theory defect
is only the universal $1/3$ per mode (a wave-function renormalization of the
modewise map, $\rho_k\to\rho_k Z_k$ with $Z_k$ unitary, cannot change class
(3), but a $k$-dependent redefinition of the mode split might); and (iii)
make contact with the Krein-space formulation of indefinite-metric QFT,
where the covariant representation lives, to phrase the representation
choice intrinsically.  We leave these to future work.

\begin{thebibliography}{9}

\bibitem{companion}
GPT-5.6.sol and A.~A.~Palm,
\emph{Canonical Positive Symplectic Diagonalization of the
Pais--Uhlenbeck Oscillator},
Paper~I of the pure-Weyl gravity programme, 2026;
file \path{paper/01-symplectic-diagonalization.pdf} in the repository
\url{https://github.com/area9innovation/weyl-gravity}.

\bibitem{companionV}
GPT-5.6.sol and A.~A.~Palm,
\emph{Interaction Obstructions, Resonant PT Breaking,
and Doubled Jordan Symmetry in Fourth-Order Theories},
Paper~V of the pure-Weyl gravity programme, 2026;
file \path{paper/05-interaction-obstructions.pdf} in the repository
\url{https://github.com/area9innovation/weyl-gravity}.

\bibitem{BM2008PRL}
C.~M.~Bender and P.~D.~Mannheim,
\emph{No-ghost theorem for the fourth-order derivative Pais--Uhlenbeck
oscillator model},
Phys.\ Rev.\ Lett.\ \textbf{100} (2008) 110402.

\bibitem{BM2008PRD}
C.~M.~Bender and P.~D.~Mannheim,
\emph{Exactly solvable PT-symmetric Hamiltonian having no Hermitian
counterpart},
Phys.\ Rev.\ D \textbf{78} (2008) 025022.

\bibitem{Bhatia2007}
R.~Bhatia, \emph{Positive Definite Matrices}, Princeton University Press,
2007,
\href{https://doi.org/10.1515/9781400827787}{doi:10.1515/9781400827787}.

\bibitem{Mostow1955}
G.~D.~Mostow,
\emph{Self-adjoint groups},
Ann.\ of Math.\ \textbf{62} (1955) 44,
\href{https://doi.org/10.2307/2007099}{doi:10.2307/2007099}.

\bibitem{KempfNess1979}
G.~Kempf and L.~Ness,
\emph{The length of vectors in representation spaces},
in: Algebraic Geometry, Lecture Notes in Math.\ \textbf{732}, Springer,
1979, 233,
\href{https://doi.org/10.1007/bfb0066647}{doi:10.1007/bfb0066647}.

\bibitem{BridsonHaefliger}
M.~R.~Bridson and A.~Haefliger,
\emph{Metric Spaces of Non-Positive Curvature},
Springer, 1999.

\bibitem{vN1939}
J.~von Neumann,
\emph{On infinite direct products},
Compos.\ Math.\ \textbf{6} (1939) 1.

\bibitem{QQG2024}
J.~Kubo and J.~Kuntz,
\emph{Unitarity through PT symmetry in quantum quadratic gravity},
arXiv:2410.08278.

\end{thebibliography}

\end{document}
